\section{Execution, recovery and payment}
\subsection{The native delegation boundary}
D1's receiver consumes a slot ID, payload and signed dispatch permit. The
installed guard requires a locally accepted allocator and checks the allocator,
holder and receiver signatures. It compares the exact receiver key, capability
subject, server/tool, input digest, request ID and capability digest with the
sealed selection. The capability must be issued by this receiving kernel and
contain one exact invoke grant, one invocation, and per-call and total monetary
ceilings no larger than the offer. The offer and slot must be live. Ordinary
native capability validation, revocation and policy checks still apply.

The installation API requires durable native admission and enables original
request retention. The receiver selects either the original argument envelope
or a governed-context layout; the caller cannot select a weaker interpretation.
In the composed layout the permit commits to every original tool argument,
including existing financial terms, without rewriting the tool payload. The guard revalidates before dispatch and checks the exact
JSON value after output transformations. A signed offer in D1 agrees to
accepted-output pricing: a qualifying predicate rejection withholds output and
uses the native reversible-hold path to produce a zero-charge denial. The
executed invocation remains consumed. This rule depends on the qualified
payment profile; a portable permit does not implement settlement by itself.

Sealing uses two durable steps: first freeze the exact selection, then persist
the signed permit before returning it. A crash between them can strand
capacity. Retrying returns the original stored permit once available, and
cannot replace its provider or request. Native custody, rather than a permit
cache, determines whether that request has already executed. Timeouts do not
authorize a replacement. Independent siblings can proceed from their own
allocations while the original operation is unresolved.

In the native fault experiment, the receiver process is killed after
\texttt{ToolReturnRecorded}. A second receiver cannot replace the original
allocation, but completes an already allocated sibling. Reopening the first
receiver recovers its recorded output without a fresh tool call. This cut has
recoverable evidence; it does not establish recovery from every possible cut.
An effect lacking authoritative completion evidence remains unknown. Ordinary
external effects acquire no exactly-once guarantee from the permit.

\subsection{Extending a running swarm}
S1 adds one operation to the existing swarm lifecycle. Given the expected
digest of the installed bundle and a signed candidate, the runtime store
checks both bundles with the ordinary live-authority verifier at store-owned
time. It then checks retention: the original graph envelope, nodes, edges,
joins, routes, witness chains, allocations and revocation epoch remain intact.
The candidate must add nodes within the original depth, fan-out and pool limits.
Each task has at most one allocation and continuation. Existing continuations
retain every field except the new graph digest and its signature. A new task
receives a new continuation; an old task cannot receive a new attempt identity.

In one immediate SQLite transaction, the installer compares the expected head,
archives the exact parent bundle and installs the successor. Competing proposals
against one parent cannot both commit. The archive binds the graph index to the
signed graph digest and full bundle digest. Admission uses the graph digest
already carried by the request, so an unfinished operation can still present
its original evidence after growth. Every fresh invocation also passes the
ordinary current capability, revocation and treaty gates.

The native kernel owns the single-use continuation resource by its stable ID.
A changed graph digest therefore cannot turn a consumed continuation into fresh
authority. Replaying a completed request resolves its original operation and
receipt. Revalidating an original physical claim reads its original graph and
treaty evidence. The installer does not need a joint transaction with every
running operation: additive growth leaves their allocations and identities
intact. Non-additive retirement would require a different authority boundary.

\subsection{One admission across three resource owners}
The composed adapter resides in the experimental funded receiver. Its protected
policy pins the allocator, swarm witness, treaty peer and runtime profile.
Opening the receiver installs the same delegation guard and runtime admission
hook before recovery. Removing that configuration changes the pinned policy;
losing its runtime database cannot silently provision empty custody.

The D1 slot ID equals the native request ID and S1 task ID. The S1 allocation ID
is the D1 allocation digest; both agree on amount and currency. The sealed
permit and native swarm witness bind the same capability. The signed route must
target this receiver. The existing F1 agreement signs the entire final request,
including the governed D1, swarm and treaty references. Its escrow allocation ID
therefore commits transitively to those references. These are different hash
domains joined by checked bindings, not interchangeable identifiers. The
agreement is downstream of the request, so no circular hash is required.

Before dispatch, the receiver checks treaty and capability authority separately
from observed escrow backing. One existing native operation acquires both the
treaty and swarm continuation resources. Its payment adapter compares the
governed payment context with the original retained intent, then binds the
original funding allocation to the operation and budget hold. A new graph
member receives neither a capability nor funds merely by appearing in a graph.

A receiver-local pre-settlement statement authenticates the retained execution
and output before a funded decision. For a federated invocation, export
revalidates the original frozen treaty, participants and admission time against
the content-addressed native outcome. It does not ask for a new remote
co-signature, imply delivery to the remote party, or settle payment. The agreed
checker and escrow decision establish the earned claim. Final bilateral
receipt delivery remains a separate operation.

The commit order is seal, bind the funded agreement, deposit backing, install
checked growth, admit and dispatch, then record acceptance. No distributed
transaction spans those owners. A crash between steps may strand unused
capacity or funds until the existing refund rules apply. It cannot silently
replace the original receiver, request or allocation.

\subsection{The funded state machine}
F1's financial state machine has seven states after funding:
\[
\textsf{Funded}\to\textsf{Submitted}\to
\begin{cases}
\textsf{Payable}\to\textsf{Paid},\\
\textsf{Rejected}\to\textsf{Refunded}.
\end{cases}
\]
A funded or submitted job without a decision can become
$\textsf{TimedOut}$ and then $\textsf{Refunded}$. There is no edge from
$\textsf{Payable}$ to timeout or refund.

Let $t_s<t_c<t_r<t_f$ denote submission, challenge, resolution and refund
boundaries. The beneficiary submits by $t_s$. A first verifier decision is
recorded only when $t_c<t\leq t_r$. Without a decision, refund is allowed only
when $t>t_f$. The interval before decision is a timing boundary, not a claim
that F1 implements a public challenge court. The exact verifier decision may
be carried by any courier. Its first valid recording determines acceptance;
an identical replay changes nothing, and a conflicting replay fails.

A claim is \emph{earned} here when an accepting decision has been recorded in
the settlement state under the selected finality rule. Merely producing an
output or obtaining an off-chain signature does not earn this right. After
that transition, the beneficiary may withdraw without a new payer signature.
The reference contract imposes no later expiry on an earned withdrawal.
Administrative pause and verifier-registry changes govern new funding and do
not revise an existing claim. The recipient still needs transaction inclusion,
fees and a functioning asset contract.

Local execution proceeds under its own durable operation identity. The kernel
commits admission before dispatch, retains observed output and produces signed
evidence. The funded integration carries the original allocation, operation,
hold and authorization identifiers through verification and recovery. Recovery
reopens that history; it must not mint a fresh admission simply because a
response was lost. A return whose occurrence cannot be established remains an
unknown terminal. A later financial resolution has its own authority and does
not rewrite the original observation.

For actual rail transactions, the research implementation records the exact
signed transaction before broadcast. After a crash it reobserves or rebroadcasts
those same bytes. A changed nonce, transaction or funding context is not an
ordinary retry. This approach handles particular durable cutpoints; it is not
an exactly-once guarantee for arbitrary tool effects.

\subsection{Successful exchange and failed intermediary}
In a successful exchange, the buyer and provider sign the terms, the buyer
funds $f_1$, and the provider admits one bound operation. The result and custody
record identify the exact submitted bytes. The verifier reruns the selected
check, records acceptance, and the provider withdraws 100. An offline reader
can verify the signed evidence without possessing either party's private keys;
current spendability still requires settlement state.

For the failed intermediary, the specialist's 60-unit child claim becomes
$\textsf{Payable}$ but remains uncollected. The parent produces no eligible
submission and its 100 returns to the buyer after timeout. The specialist then
withdraws the original 60. The provider bears the loss because it funded the
child. This is neither a transfer of the buyer's risk to an unpaid specialist
nor a claim that failed computation is free.

If instead the verifier never records a child decision, the child may time out
and refund its payer despite the specialist having performed work. F1 protects
recorded earned claims; it cannot guarantee timely adjudication against its own
failed verifier. That distinction is the decisive difference between financial
safety and an unconditional promise to pay for every useful result.
